\section{Further research questions and proposed priorities}
\label{sec:research}

The arithmetic improvement is useful enough to integrate as an optional
capability. The evidence suggests a different order of priorities for obtaining
practical speed-ups and for proving a quasi-polynomial theorem. The former
needs better scheduling and reuse; the latter needs a representation theorem
and a bound on decisive work. The following questions state both the proposed
work and the evidence that would count as progress.

\subsection{Near-term implementation and experimental questions}

\paragraph{Q1. When should a suffix pay for a response?}
Can a scheduler predict the number of distinct useful matching and partition
queries before the suffix changes? A candidate decision rule should compare
measured direct-query cost with preparation cost plus expected boundary-query
cost, using $v-t$, $t$, projected query count, and the number of likely prime
passes. It should also account for the chance that the exact Euler prefilter
or first few direct queries stop the scan. A concrete milestone is an opt-in
hybrid that selects direct integer, rational response, or modular response
on a held-out corpus and improves total recognition time with every failed
probe charged. The current seven raw regressions make this the first
performance priority. A model fitted only to the favorable pure-tail stream
would not be adequate evidence.

\paragraph{Q2. Can a response survive a change of suffix?}
Removing one crossing changes both edges and the cut-face fragmentation. It can
split a region, create new terminals, and change the radical dimension. A
correct incremental update therefore needs a precise map between successive
vertex spaces, not just a rank-one edge update to a fixed Laplacian. Develop
an update calculus for fragment splits, terminal promotions, and edge deletion,
with a rebuild path when the pivot profile changes. The useful theorem would
bound amortized update work in terms of the changed frontier neighborhood.
The useful experiment would include many consecutive observed stages with
few queries per stage, where full rebuilding currently loses.

\paragraph{Q3. How should articulation blocks and sparse elimination combine?}
The maintained direct observer already factors at articulation points. A
boundary response could eliminate terminal-free blocks separately and retain
small interfaces at blocks that contain frontier terminals. The challenge is
to prove the gluing formula with signed weights, zero tree factors, and singular
interiors; a division by an apparently harmless block determinant can fail
modulo a prime. Establish a singular-safe block response and compare it with
both existing articulation factoring and dense modular preparation. Measure
fill-in, temporary memory, and cost at exceptional primes, as well as elapsed
time on nonsingular examples.

\paragraph{Q4. Which certified coefficient bounds repay their cost?}
The bound $C_a\le s_a2^{n-1}$ is universal but can be loose. Candidate refinements
include the unsigned spanning-tree count of the actual graph, products of
blockwise bounds, bridge reductions, and degree-based determinant bounds. A
bound should be computed once and reused across completions when possible;
otherwise its cost can exceed the modular observation it was meant to save.
The target is a proved $C_a'$ together with a measured reduction in prime passes
or exact-threshold completion time. Signed cancellation may improve the actual
coefficient, but cannot justify shrinking a certified bound without evidence.

\paragraph{Q5. Can prime scheduling exploit the exposure theorem?}
\Cref{prop:prime-exposure} bounds how many primes can hide an existing
obstruction. Use it to design a policy choosing between one more independent
prime, another CRT pass, and continuing the scan. An asymptotic version should
supply enough primes automatically and charge their generation and validation.
A practical version should distinguish a large default prime from a random
prime-pool experiment, retain monotone CRT information, and report the cost
of inconclusive passes. Prime choice must never change the meaning of a
positive result. The 70 misses at prime three are a useful adversarial check,
not a reason to label the method probabilistic or unsound.

\paragraph{Q6. What is the true memory cost of continuation reuse?}
Only the active residual matrices are boundary-sized. Completed-value keys and
records can accumulate across stages, and algebraic generator storage may
dominate everything. Introduce a memory ledger for geometry, prime responses,
partition caches, matching keys, exact Euler integers, and scanner objects.
Study bounded caches or stage-local eviction with reconstruction charged.
A credible performance report should include peak memory and the number of
retained objects, not only the largest query matrix. This may reveal a regime
where recomputation is preferable to retaining all exact completion values.

\subsection{Questions that address the quasi-polynomial barrier}

\paragraph{Q7. Which continuation behaviors admit a succinct quotient?}
Equality of boundary matchings is too coarse to identify chain complexes with
different differentials. Seek an equivalence relation on relative complexes
that preserves the needed threshold decision under all admissible future
crossings or actual suffixes, and that is stable under tensoring, delooping,
and cancellation. The essential result would bound the number and bit size of
its states by $2^{(\log N)^{O(1)}}$ and give constructive transition rules.
A finite table of scalar Euler observations alone is unlikely to determine
all future homological behavior; this should be tested with counterexamples
before adopting it as a state representation. Any successful quotient must
retain enough information to justify the marked homotopies, not merely $d^2=0$.

\paragraph{Q8. Is there a family theorem for the first decisive observation?}
For a specified diagram family and order policy, can one bound the number of
represented objects and stages before $B_4>1$, or before a rank-two contraction
is exhibited? The arithmetic completion theorem starts only after an exact
obstruction is present. It says nothing about when that happens. A useful
intermediate target is an explicit family beyond the already solved structural
cases, with a proved stopping-stage bound and a certified complexity expression
in its parameters. Failed examples should retain the exact sequence of prefix
bounds, so a proposed detection lemma can be falsified by genuine complexes.

\paragraph{Q9. Can larger residue systems remain efficiently observable?}
Residues modulo eight or higher powers of two refine quantum grading and can
retain cancellations lost modulo four. Investigate whether their evaluations
admit a manageable transfer representation on bounded interfaces. The special
value at $i$ has a spanning-tree determinant formula; higher roots do not
inherit that formula automatically. A successful extension needs three proofs:
that the residue norm still bounds completed homology, that the relevant
integer or cyclotomic data are computed exactly, and that their observation
cost has the claimed dependence on boundary width. A larger invariant that
requires exponential work in the suffix would not resolve the present barrier.

\paragraph{Q10. Can homological information be compressed without changing fields?}
The auxiliary primes in this paper reduce integer Euler data. They cannot be
used to compute a generic odd-characteristic Khovanov rank and silently treat
it as the rank over $\F_2$. Investigate compact descriptions of persistent
cancellations, module decompositions, or continuation-compatible rank defects
\emph{over the actual coefficient field}. The maintained interval and scalar-
splitting backends provide concrete comparison points. A new theorem should
state exactly which operations preserve the compressed description and what
happens when that description must expand. The three-braid structural formulas
suggest that family-specific symbolic answers can be much smaller than
explicit bases \cite{kelomaki}.

\paragraph{Q11. What geometric decomposition gives the right width measure?}
A polylogarithmic number of frontier darts would bound the oracle table, but
it is not guaranteed by the current crossing-order search. Investigate
hierarchical decompositions whose interfaces carry both the tangle data and
sufficient algebraic summaries. A useful theorem must control interface size,
number of pieces, transition cost, and the growth of multiplicities and
coefficient descriptions simultaneously. If only a planar separator estimate
with a much larger boundary is available, state the resulting subexponential
bound honestly; it does not become quasi-polynomial through terminology.
Connections to hierarchy-based topological methods are promising only after
the translation cost is made explicit.

\paragraph{Q12. Can independently verified compressed contractions replace replay?}
The current verifier reconstructs the entire prefix. A more useful certificate
would encode contractions and direct-sum decompositions in a form whose local
identities can be checked without expanding every intermediate complex. Study
straight-line descriptions of maps and homotopies, canonical sharing, and
certified transfer across a decomposition tree. The certificate must include
its own size bound and account for zero tests and equality of compressed maps.
A short list of moves is insufficient if one move refers to an exponentially
large object. A constructive bound here would address both practical replay
cost and a major part of \cref{prop:conditional}.

\subsection{External algorithm audits and integration milestones}

Two literature tracks deserve separate, explicit audits. Lackenby's hierarchy
program supplies a geometric framework and efficiently checkable objects;
an implementation study should identify which construction and encoding bounds
are already proved and which speed-up lemmas remain proposed
\cite{lackenby2026,lackenbytalk}. Musick's September 2026 polynomial-time claim
should be assessed independently: formalize its state representation, verify
that every claimed local move is discoverable within the compressed encoding,
and prove a bound on the total descriptions created, including overwritten
or retained grammar versions \cite{musick}. This article neither adopts that
claim nor supplies a counterexample. Its relevance is that a serious current
research program should investigate a directly competing claimed algorithm
instead of declaring it absent.

For ProveIt, the immediate integration sequence is concrete. Add the optional
backend and its replay tests against the pinned source. Preserve the recorded
arithmetic and actual-diagram audits. Keep the raw timing samples and the
negative complete-scan results next to the favorable response experiment.
Then evaluate a scheduling change as a separate patch with a new held-out
corpus. The inherited normal-surface I/O defect should also receive a separate
focused repair rather than an unexplained test exclusion.

The next mathematical milestone should be a family theorem controlling the
size of a continuation-compatible representation through a decisive event.
The modular response provides a proved inexpensive observation interface to
attach to that theorem. Establishing the representation and stopping bounds
is the work that would turn the present performance improvement into progress
on a general quasi-polynomial recognition guarantee.
